\section{联合调度约束}
\label{sec:constraints}

\subsection{设备可行域}

联合调度的可行域由4.2--4.7的设备方程和公共母线平衡共同构成。源侧满足可用功率、
状态和爬坡约束；电池满足能量递推、SOC和构网备用约束；制氢系统满足模块启停、爬坡、
产氢和库存约束；算力系统满足任务队列、设施容量、功率和SLA约束；产品结算量受海缆、
储运、光纤和终端需求限制。

设备故障、检修和极端天气通过时变可用率进入容量上界：
\begin{equation}
0\le x_{j,t}\le A_{j,t}^{evt}\,\overline x_{j,t},
\qquad j\in\mathcal J^{asset},
\label{eq:availability-general}
\end{equation}
其中$A_{j,t}^{evt}\in[0,1]$。某条产品路径关闭时，仅将该路径的专属功率、流量和收入变量
置零，公共电源、储能和内部关键负荷仍保留在模型中。

\subsection{产品分配计划约束}

为表达在线策略给出的电力外送、制氢和柔性算力分配计划，定义
\begin{equation}
p_{E,t}=P_t^{cab,s},\qquad
p_{H,t}=P_t^{EL},\qquad
p_{C,t}=P_t^{DC,flex},\qquad
p_{\Sigma,t}=p_{E,t}+p_{H,t}+p_{C,t}.
\label{eq:hourly-mix-def}
\end{equation}
其中$P_t^{DC,flex}$表示扣除基础设施功率和刚性任务功率后的柔性算力用电。给定计划份额
$q_{E,t}+q_{H,t}+q_{C,t}=1$及允许偏差$\delta$，对$i\in\{E,H,C\}$有
\begin{equation}
\max(0,q_{i,t}-\delta)p_{\Sigma,t}
\le p_{i,t}\le
\min(1,q_{i,t}+\delta)p_{\Sigma,t}.
\label{eq:hourly-mix}
\end{equation}
该约束只作用于当期实际可分配功率。若$p_{\Sigma,t}=0$，各产品分配量同时为零；计划份额
不用于反推设备装机比例，也不要求在资源不足时强制维持固定比例。

\subsection{服务等级、终端状态与极端事件}

内部关键负荷、刚性海洋用能和刚性算力任务分别设置未供变量。联合模型可以通过高惩罚
体现服务优先级，也可以在$\varepsilon$约束中限制未供能上限。惩罚系数用于区分服务等级，
不改变设备的物理容量。

有限时域仿真可采用循环终端或最低库存终端：
\begin{equation}
E_T^B=E_0^B\ \text{或}\ E_T^B\ge E^{B,target},\qquad
H_T=H_0\ \text{或}\ H_T\ge H^{target}.
\label{eq:terminal-rules}
\end{equation}
循环终端适用于典型周期比较；最低库存终端适用于安全储备场景；年度逐时运行则连续传递
状态，并通过滚动策略控制库存水平。

极端事件状态定义为
\begin{equation}
z_t^{evt}\in\{\mathrm{NORMAL},\mathrm{WARNING},
\mathrm{PASSAGE},\mathrm{RECOVERY}\}.
\label{eq:event-state}
\end{equation}
预警阶段可降低可选外送、制氢和柔性算力，优先恢复电池SOC和必要库存；过境阶段根据设备
保护状态和通信条件限制运行；恢复阶段依据实际可用率和启动约束逐步解除限制。各阶段的
降额幅度、储备目标和恢复速度由气象预报、设备规范和应急预案确定。

\subsection{变量域与模型线性化}

功率、电量、产量、库存、任务完成量和未供量为连续变量；电池充放状态、电解槽在线模块数、
设备启动和可选氢转电状态为整数或二元变量。采用固定或分段线性的SEC、PUE和海缆损耗时，
联合运行问题可写成MILP。

第五章年度基准仿真在数据不足时采用外生逐时PUE和固定海缆损耗率；该处理属于场景简化。
当设备曲线和任务数据完善后，可替换为4.3与4.6中的分段线性模型，并在求解后使用原始物理
函数复算，以检查近似误差。设备容量在本章作为已知参数，容量--运行联合优化需在独立的
规划模型中处理。
